\chapter{Conclusions}

Conclusions










%
%A method can be characterized by its input, techniques, and output.
%The input of our work is diff files, the techniques are neural networks, and the output is classification.
%Thus, in this section we first describe the format of diff files, and then introduce the related work in the aspects of input, techniques, and output.
%
%
%
%
%Code summarization is a sub-field of natural language processing (NLP),
%here the input is the program source code and the output is the summary.
%A large body of work has been carried out in this field with assorted techniques.
%Despite the fact that code summarization does not match our work in all the three aspects,
%some observations inspired us.
%
%
%
%When the input is source code rather than diff,
%the field becomes program analysis, where more and more work has employed neural network techniques.
%\etal{Alexandru}~\cite{alexandru2017replicating} used NMT to annotate source code tokens with typing,
%and they also implemented an {\sc antlr}-based parser.
%\textsc{JSNice}~\cite{raychev2015predicting} is a neural network model that predicts names of JavaScript identifiers and type annotations of variables.
%\etal{Mou}~\cite{mou2016convolutional} used a convolutional neural network to classify programs by functionalities, such as string manipulations and matrix operations.
%Other uses of neural networks in program analysis include
%detection of variable misuse~\cite{morgachev2019detection},
%bug localization~\cite{huo2016learning},
%API suggestion~\cite{gu2016deep}, and
%code completion~\cite{raychev2014code}.
%
%Taking diff as input, program analysis is specialized to diff analysis.
%\etal{Moreno}~\cite{moreno2016arena} generated release notes by reading changes to program source code and documentations, as well as taking into consideration issues from software repositories.
%ChangeScribe~\cite{linares2015changescribe},
%ChangeDoc~\cite{huang2020learning}, and
%DeltaDoc~\cite{buse2010automatically} read diff and followed a set of rules to generate commit messages.
%None of them used machine learning, and howbeit they all processed diff,
%their goal was not to classify the diff.
%They do not accept the {\gitdiff} format which shows barely changed lines of changed files.
%
%
%
%
%
%
%
%
%
%There are existing works in the field of program analysis,
%but research is struggling to answer questions like ``Why does this method crash?'', ``Is this code thread-safe?''~\cite{alexandru2017replicating}.
%To enrich the experience towards solving the broader problem of code classification,
%we carry out this experiment to tackle one specific type of code modification recognition
%--- renaming of identifiers.
%In this work, we propose a probabilistic machine-learning classification model,
%per taxonomy of~\cite{allamanis2018survey},
%to automatically identify the commits to a version-control system that only consist of identifier renaming.
%This technique is novel in that the neural network is able to read merely the diff and give a confidence value of whether it belongs to the category.
%On the contrary, the traditional way of detecting such changes is done by analyzing the full source code before and after the commit,
%which is less efficient and often complicated with partial source code.
%For typical software projects nowadays each with a lot of imported libraries,
%the ability to carry out analysis based on local partial source code is certainly an advantage.
%
%Our model can be used as a component in the Continuous Integration (CI) with a Pull-Request (PR) workflow, for instance, on GitHub Actions.
%When a developer submits a pull request, CI runs to build and test the code and can add labels to the PR.
%Our model plays a role in CI that we take the unified diff of the PR as input,
%and classify whether the diff is identifier renaming, ``yes'' or ``no''.
%Then, CI can label the PR or assign it to the corresponding code reviewers.
%For our model to work, if one commit changes multiple files,
%the changes to these files should be concatenated to a single diff file,
%as the~{\gitdiff} program does.
%
%Our work focuses on the classification of diff in a typical programming language
%--- Java diff files,
%because Java is a popular language and the syntax is relatively simple.
%We specifically target Java 7, since Java 8 and beyond have introduced many new syntactical improvements,
%which could cause complication in the parsing and needlessly shift our focus away.
%
%There is no off-the-shelf dataset of commits and labels of whether the commit is identifier renaming.
%Although \etal{Jiang} claimed their dataset contained only Java source code,
%we found it actually containing many non-Java contents.
%We filtered out non-Java code and 62K commits remain.
%To further expand the dataset, we employed the crawling GitHub tool provided by~\cite{alexandru2017replicating} to collect additional commits from top Java repos on GitHub in January 2021.
%The final dataset obtained has \num{73080} examples.
%
%We contribute in this work a probabilistic neural network classification model with diff files as input, telling whether the input is renaming or not.
%The model was specifically trained against the Java programming language, and
%we believe minor tweaks can allow it to work on other programming languages.
%Besides the neural network model, we also present a dataset that labels 73K Java diff files with whether they are renaming commits,
%together with a syntax analyzer for Java diff files based on {\sc antlr}.
%
%
%
%
%
%
%\begin{equation}\label{eq:sample1}
%	a^2 + b^2 = c^2
%\end{equation}
%
%\begin{equation}\label{eq:sample2}
%	\begin{dcases}
%		a + b = 2\\
%		a - b = 0
%	\end{dcases}
%	\implies
%	\begin{dcases}
%		a = 1\\
%		b = 1
%	\end{dcases}
%\end{equation}
%
%\autoref{eq:sample1} and \ref{eq:sample2} as sample equation.
%
%
%\begin{table} [h]
%	\centering
%	\caption{Sample of table.}
%	\begin{tabular}{|c|c|c|c|}
%		\hline
%		\textbf{A} & \textbf{B} & \textbf{C} & \textbf{D} \\
%		\hline
%		1 & 2 & 3 & 4 \\
%		\hline
%		5 & 6 & 7 & 8 \\
%		\hline
%	\end{tabular}
%	\label{tab:sample}
%\end{table}
%
%\autoref{tab:sample} as sample table.
%
%\begin{figure} [h]
%	\centering
%	\includegraphics[width=0.5\linewidth]{MPI.pdf}
%	\caption{MPI logo}
%	\label{fig:sample}
%\end{figure}
%
%\autoref{fig:sample} as sample figure.